\chapter{Linear Algebra: Stuck in the Matrix}

Up to this point we have written equations independent of each other.

Earlier we described our vectors by
\begin{equation}
    \vb{r} = (r_x,r_y,r_z) = r_x \vu{x} + r_y \vu{y} + r_z \vu{z} .
\end{equation}
However, there is another way.  We can express a vector as a either a single row
or single column matrix, e.g.
\begin{align}
    \mat{r} & = \begin{bmatrix}
            r_x & r_y & r_z
        \end{bmatrix} && \text{row vector} , \\
    & = \begin{bmatrix}
            r_x \\ r_y \\ r_z
        \end{bmatrix} && \text{column vector} . \\
\end{align}
Note: we change the notation for a vector from $\vb{r}$ to $\mat{r}$ when
we express it as a matrix.  While this change in notation is not really
necessary it is a commonly used convention.


\section{Matrix Math}

Matrix notation allows us to compactly represent systems of linear equations.
For example, a system of three equations in three unknowns can be written as
\begin{align}
    a_{11} x + a_{12} y + a_{13} z &= b_1 ,\\
    a_{21} x + a_{22} y + a_{23} z &= b_2 ,\\
    a_{31} x + a_{32} y + a_{33} z &= b_3 .
\end{align}
This can be expressed in matrix form as
\begin{equation}
    \mathbf{A} \vb{x} = \vb{b},
\end{equation}
where
\begin{equation}
    \mathbf{A} =
    \begin{bmatrix}
        a_{11} & a_{12} & a_{13} \\
        a_{21} & a_{22} & a_{23} \\
        a_{31} & a_{32} & a_{33}
    \end{bmatrix}, \quad
    \vb{x} =
    \begin{bmatrix}
        x \\ y \\ z
    \end{bmatrix}, \quad
    \vb{b} =
    \begin{bmatrix}
        b_1 \\ b_2 \\ b_3
    \end{bmatrix}.
\end{equation}

Matrix multiplication follows a row-by-column rule. The $(i,j)$ element of the
product $\mathbf{C} = \mathbf{A}\mathbf{B}$ is given by
\begin{equation}
    c_{ij} = \sum_{k} a_{ik} b_{kj}.
\end{equation}
This operation is only defined when the number of columns in $\mathbf{A}$ equals
the number of rows in $\mathbf{B}$.

Note that matrix multiplication is generally \emph{not commutative}, i.e.
\begin{equation}
    \mathbf{A}\mathbf{B} \neq \mathbf{B}\mathbf{A}.
\end{equation}


\section{Matrix Operators}

\subsection{Determinant}

The determinant is a scalar quantity associated with a square matrix. For a
$2\times2$ matrix,
\begin{equation}
    \mathbf{A} =
    \begin{bmatrix}
        a & b \\
        c & d
    \end{bmatrix},
\end{equation}
the determinant is
\begin{equation}
    \det(\mathbf{A}) = ad - bc.
\end{equation}

For a $3\times3$ matrix,
\begin{equation}
    \det(\mathbf{A}) =
    \begin{vmatrix}
        a_{11} & a_{12} & a_{13} \\
        a_{21} & a_{22} & a_{23} \\
        a_{31} & a_{32} & a_{33}
    \end{vmatrix},
\end{equation}
which can be expanded along the first row as
\begin{equation}
    \det(\mathbf{A}) =
    a_{11}
    \begin{vmatrix}
        a_{22} & a_{23} \\
        a_{32} & a_{33}
    \end{vmatrix}
    - a_{12}
    \begin{vmatrix}
        a_{21} & a_{23} \\
        a_{31} & a_{33}
    \end{vmatrix}
    + a_{13}
    \begin{vmatrix}
        a_{21} & a_{22} \\
        a_{31} & a_{32}
    \end{vmatrix}.
\end{equation}

The determinant has several important interpretations:
\begin{itemize}
    \item It measures the scaling of volume under the linear transformation.
    \item If $\det(\mathbf{A}) = 0$, the matrix is singular (non-invertible).
    \item The sign indicates orientation (e.g.\ right-handed vs left-handed).
\end{itemize}

\subsection{Transpose}

The transpose of a matrix, denoted $\mathbf{A}^T$, is obtained by swapping rows
and columns:
\begin{equation}
    (\mathbf{A}^T)_{ij} = a_{ji}.
\end{equation}

For example,
\begin{equation}
    \mathbf{A} =
    \begin{bmatrix}
        a_{11} & a_{12} \\
        a_{21} & a_{22}
    \end{bmatrix}
    \quad \Rightarrow \quad
    \mathbf{A}^T =
    \begin{bmatrix}
        a_{11} & a_{21} \\
        a_{12} & a_{22}
    \end{bmatrix}.
\end{equation}

Important properties include
\begin{align}
    (\mathbf{A}^T)^T &= \mathbf{A}, \\
    (\mathbf{A}\mathbf{B})^T &= \mathbf{B}^T \mathbf{A}^T.
\end{align}

A matrix is called \emph{symmetric} if $\mathbf{A}^T = \mathbf{A}$.

\subsection{Inverse}

The inverse of a matrix $\mathbf{A}$, denoted $\mathbf{A}^{-1}$, satisfies
\begin{equation}
    \mathbf{A}^{-1} \mathbf{A} = \mathbf{A} \mathbf{A}^{-1} = \mathbf{I},
\end{equation}
where $\mathbf{I}$ is the identity matrix.

For a $2\times2$ matrix,
\begin{equation}
    \mathbf{A} =
    \begin{bmatrix}
        a & b \\
        c & d
    \end{bmatrix},
\end{equation}
the inverse is
\begin{equation}
    \mathbf{A}^{-1}
    =
    \frac{1}{ad - bc}
    \begin{bmatrix}
        d & -b \\
        -c & a
    \end{bmatrix},
\end{equation}
provided that $\det(\mathbf{A}) \neq 0$.

For larger matrices, the inverse is typically computed using row reduction or
numerical methods.

In practice (especially in geophysics), it is often preferable to solve
\begin{equation}
    \mathbf{A}\vb{x} = \vb{b}
\end{equation}
without explicitly computing $\mathbf{A}^{-1}$, as direct inversion can be
numerically unstable and computationally expensive.

\section{Eigenvalues and Eigenvectors}

Eigenvalues and eigenvectors describe how a matrix transforms a vector. In
particular, they identify directions that are stretched or compressed without
being rotated.

An eigenvector $\vb{v}$ of a matrix $\mathbf{A}$ satisfies
\begin{equation}
    \mathbf{A} \vb{v} = \lambda \vb{v},
\end{equation}
where $\lambda$ is the corresponding eigenvalue.

This equation says that applying $\mathbf{A}$ to $\vb{v}$ simply scales the
vector by a factor $\lambda$, without changing its direction.

To compute eigenvalues, we rewrite the equation as
\begin{equation}
    (\mathbf{A} - \lambda \mathbf{I}) \vb{v} = \vb{0}.
\end{equation}
For non-trivial solutions ($\vb{v} \neq \vb{0}$), we require
\begin{equation}
    \det(\mathbf{A} - \lambda \mathbf{I}) = 0.
\end{equation}
This is called the \emph{characteristic equation}, and its solutions give the
eigenvalues $\lambda$.

Once the eigenvalues are known, the corresponding eigenvectors are found by
solving
\begin{equation}
    (\mathbf{A} - \lambda \mathbf{I}) \vb{v} = \vb{0}.
\end{equation}

Eigenvalues and eigenvectors have important physical interpretations. In
geophysics, for example:
\begin{itemize}
    \item Eigenvectors often represent principal directions (e.g.\ principal
    stress or strain axes).
    \item Eigenvalues represent magnitudes along those directions.
\end{itemize}

For symmetric matrices, which commonly arise in physical problems, eigenvectors
are orthogonal and eigenvalues are real. This greatly simplifies their
interpretation.

\section{Singular Value Decomposition (SVD)}

The singular value decomposition (SVD) is one of the most powerful and widely
used tools in linear algebra, particularly for solving inverse problems and
analyzing ill-conditioned systems.

Any $m \times n$ matrix $\mathbf{A}$ can be written as
\begin{equation}
    \mathbf{A} = \mathbf{U} \boldsymbol{\Sigma} \mathbf{V}^T,
\end{equation}
where
\begin{itemize}
    \item $\mathbf{U}$ is an $m \times m$ orthogonal matrix,
    \item $\mathbf{V}$ is an $n \times n$ orthogonal matrix,
    \item $\boldsymbol{\Sigma}$ is an $m \times n$ diagonal matrix containing
    non-negative values.
\end{itemize}

The diagonal entries of $\boldsymbol{\Sigma}$,
\begin{equation}
    \sigma_1 \geq \sigma_2 \geq \cdots \geq 0,
\end{equation}
are called the \emph{singular values} of $\mathbf{A}$.

\subsection*{Interpretation}

The SVD provides a geometric interpretation of a matrix as a sequence of three
operations:
\begin{enumerate}
    \item $\mathbf{V}^T$: rotates the input vector into a new coordinate system,
    \item $\boldsymbol{\Sigma}$: stretches (or compresses) along orthogonal
    directions,
    \item $\mathbf{U}$: rotates the result into the final coordinate system.
\end{enumerate}

Thus, $\mathbf{A}$ maps a unit sphere into an ellipsoid, whose principal axes
are given by the columns of $\mathbf{U}$ and whose lengths are the singular
values $\sigma_i$.

\subsection*{Connection to Eigenvalues}

The singular values are related to the eigenvalues of $\mathbf{A}^T \mathbf{A}$:
\begin{equation}
    \mathbf{A}^T \mathbf{A} \vb{v}_i = \sigma_i^2 \vb{v}_i,
\end{equation}
where $\vb{v}_i$ are the columns of $\mathbf{V}$. Similarly,
\begin{equation}
    \mathbf{A} \mathbf{A}^T \vb{u}_i = \sigma_i^2 \vb{u}_i,
\end{equation}
where $\vb{u}_i$ are the columns of $\mathbf{U}$.

\subsection*{Rank and Singular Values}

The rank of $\mathbf{A}$ is equal to the number of non-zero singular values:
\begin{equation}
    \mathrm{rank}(\mathbf{A}) = \# \{ \sigma_i > 0 \}.
\end{equation}

Small singular values indicate directions in which the matrix has little
sensitivity. In practice, very small $\sigma_i$ are often treated as zero,
revealing effective rank deficiency.

\subsection*{Solving Linear Systems}

The SVD provides a stable way to solve
\begin{equation}
    \mathbf{A} \vb{x} = \vb{b},
\end{equation}
especially when $\mathbf{A}$ is ill-conditioned or not square.

The solution can be written as
\begin{equation}
    \vb{x} = \mathbf{V} \boldsymbol{\Sigma}^{-1} \mathbf{U}^T \vb{b},
\end{equation}
where $\boldsymbol{\Sigma}^{-1}$ is formed by inverting the non-zero singular
values.

If some singular values are zero (or very small), the inverse does not exist in
the usual sense. In this case, we use the \emph{pseudo-inverse},
\begin{equation}
    \mathbf{A}^{+} = \mathbf{V} \boldsymbol{\Sigma}^{+} \mathbf{U}^T,
\end{equation}
which provides the least-squares solution.

\subsection*{Importance in Geophysics}

SVD is particularly useful in geophysical applications:
\begin{itemize}
    \item It reveals rank deficiency and non-uniqueness in inverse problems.
    \item It identifies poorly constrained model parameters (small $\sigma_i$).
    \item It provides a natural framework for regularization by damping small
    singular values.
\end{itemize}

For these reasons, SVD is a central tool in modern inversion theory.

\section{Matrix Rank}

The \emph{rank} of a matrix is the number of linearly independent rows or
columns in the matrix. It provides a measure of how much independent
information the matrix contains.

For a matrix $\mathbf{A}$,
\begin{equation}
    \mathrm{rank}(\mathbf{A}) = \text{number of linearly independent rows}
    = \text{number of linearly independent columns}.
\end{equation}

Some important consequences:
\begin{itemize}
    \item If $\mathbf{A}$ is an $n \times n$ matrix and
    $\mathrm{rank}(\mathbf{A}) = n$, then $\mathbf{A}$ is full rank and
    invertible.
    \item If $\mathrm{rank}(\mathbf{A}) < n$, then $\mathbf{A}$ is singular and
    does not have an inverse.
    \item The rank determines the number of independent equations in a system
    $\mathbf{A}\vb{x} = \vb{b}$.
\end{itemize}

Geometrically, the rank tells us the dimension of the space spanned by the
matrix columns:
\begin{itemize}
    \item Rank 1: all columns lie along a line.
    \item Rank 2: columns span a plane.
    \item Rank 3: columns span a volume (in $\mathbb{R}^3$).
\end{itemize}

In practice, rank is determined by performing row reduction to bring the matrix
into row echelon form. The number of non-zero rows (or pivots) is the rank.

Rank is particularly important in geophysical inversion problems:
\begin{itemize}
    \item A rank-deficient system indicates redundancy or insufficient data.
    \item It may lead to non-unique solutions.
    \item Regularization methods are often used to handle such cases.
\end{itemize}